\section{Capacity to Ship Packages within D Days}
\label{sec:bs-ship-packages}

\problemstatement{We are given an array $\textit{weights}$ of $n$ package
weights, which must be shipped \textbf{in the given order} (no
reordering). Each day, the ship loads a contiguous run of packages, in
order, whose total weight does not exceed the ship's capacity, then sails;
the next day starts loading from the next unshipped package. We must find
the \textbf{minimum ship capacity} that allows all packages to be shipped
within $D$ days.}

\begin{patternbox}
\emph{``Minimum capacity such that all items can be processed within a day
limit, items must stay in order''} is another binary-search-on-the-answer
problem, but the feasibility check itself is structurally different from
Sections~\ref{sec:bs-koko}--\ref{sec:bs-smallest-divisor}: rather than
summing independent per-item costs, we must \textbf{greedily simulate} a
day-by-day loading process, because the order constraint means packages
cannot be considered independently of their neighbors. Whenever a
binary-search-on-the-answer problem also mentions \emph{order} or
\emph{contiguity} constraints on how items are grouped, expect the
feasibility check itself to require a small greedy simulation rather than
a single summation formula.
\end{patternbox}

\subsection*{Understanding the problem carefully}

For a fixed candidate capacity $c$, simulate the shipping process greedily:
start a new day, and keep adding the next package's weight to today's load
as long as it does not exceed $c$; the moment adding the next package
would exceed $c$, end the day and start a new one with that package. This
greedy simulation is itself optimal for \emph{counting the minimum days
needed for a fixed capacity} -- packing each day as full as possible before
moving to the next day can never use \emph{more} days than any other valid
loading strategy, because deliberately under-loading a day only defers
weight to later days without ever reducing the total weight that must
still be shipped.

Define $\textit{daysNeeded}(c)$ as the number of days this greedy
simulation takes. As $c$ increases, each day can carry at least as much as
before, so $\textit{daysNeeded}(c)$ is non-increasing in $c$.

\begin{ideabox}[Candidate Space and Predicate]
The candidate space is integer capacities in
$[\max(\textit{weights}), \sum(\textit{weights})]$: the capacity must be
at least the heaviest single package (otherwise that package alone could
never be shipped), and capacity equal to the total weight always finishes
in exactly $1$ day. The predicate is
$P(c) = (\textit{daysNeeded}(c) \le D)$, and we want the smallest $c$
where $P$ holds.
\end{ideabox}

\subsection*{Why the greedy simulation inside the predicate is valid}

This is a small but important instance of \emph{nesting} a greedy
algorithm inside a binary search predicate: the binary search handles
\emph{which} capacity to test, while a greedy simulation -- packing each
day as full as the fixed capacity allows before moving on -- handles
\emph{whether} that capacity is feasible. The greedy simulation's
correctness rests on the exchange argument that ending a day early
(loading less than the maximum the capacity allows) can never help and can
only push more weight into later days, since the in-order constraint means
no future package can be moved earlier to compensate. This nested
structure -- binary search over a monotonic outer property, with a greedy
or simulation-based feasibility check inside -- is common throughout the
``binary search on the answer'' family and is exactly what makes capacity
$c$ feasible or infeasible a well-defined, monotonic question.

\subsection*{Algorithm}

\begin{enumerate}
  \item Initialize $\textit{lo} \gets \max(\textit{weights})$,
        $\textit{hi} \gets \sum(\textit{weights})$.
  \item While $\textit{lo} \le \textit{hi}$:
  \begin{enumerate}
    \item $\textit{mid} \gets \textit{lo}+(\textit{hi}-\textit{lo})/2$.
    \item Compute $\textit{daysNeeded}(\textit{mid})$ by greedily
          simulating loading days.
    \item If $\textit{daysNeeded}(\textit{mid}) \le D$: feasible; try
          smaller: $\textit{hi} \gets \textit{mid}-1$.
    \item Else: infeasible; try larger: $\textit{lo} \gets \textit{mid}+1$.
  \end{enumerate}
  \item Return \textit{lo}.
\end{enumerate}

\begin{algorithm}[H]
\caption{Minimum Capacity to Ship within D Days}
\begin{algorithmic}[1]
\Procedure{ShipWithinDays}{\textit{weights}, $D$}
  \State $\textit{lo} \gets \max(\textit{weights})$, $\textit{hi} \gets \sum(\textit{weights})$
  \While{$\textit{lo} \le \textit{hi}$}
    \State $\textit{mid} \gets \textit{lo}+(\textit{hi}-\textit{lo})/2$
    \State $\textit{days} \gets \Call{DaysNeeded}{\textit{weights}, \textit{mid}}$
    \If{$\textit{days} \le D$}
      \State $\textit{hi} \gets \textit{mid} - 1$
    \Else
      \State $\textit{lo} \gets \textit{mid} + 1$
    \EndIf
  \EndWhile
  \State \Return \textit{lo}
\EndProcedure
\end{algorithmic}
\end{algorithm}

\subsection*{Dry run}

\dryrun{Take $\textit{weights} = [1,2,3,4,5,6,7,8,9,10]$, $D = 5$.}

$\max = 10$, $\sum = 55$.

\begin{itemize}
  \item $\textit{lo}=10,\textit{hi}=55$: $\textit{mid}=32$. Simulating: day
        loads accumulate $1{+}2{+}3{+}4{+}5{+}6{+}7=28$ (next, $8$, would
        make $36>32$, so day ends); day 2 loads $8{+}9{+}10=27$. Total
        $\textit{daysNeeded}=2 \le 5$. Feasible. $\textit{hi}=31$.
  \item Binary search continues narrowing downward; eventually it
        converges to $\textit{lo}=15$, the minimum feasible capacity (with
        capacity $15$: day loads $1{+}2{+}3{+}4{+}5=15$; day 2 loads
        $6{+}7=13$ (next, $8$, would make $21>15$); day 3 loads $8$ alone
        (next, $9$, would make $17>15$); day 4 loads $9$ alone (next,
        $10$, would make $19>15$); day 5 loads $10$. Total
        $\textit{daysNeeded}=5 \le 5$ -- exactly feasible, and capacity
        $14$ can be checked to need $6$ days, confirming $15$ is minimal).
\end{itemize}

The answer is $15$.

\subsection*{C++ implementation}

\begin{lstlisting}
#include <bits/stdc++.h>
using namespace std;

int daysNeeded(vector<int>& weights, int capacity) {
    int days = 1;
    long long currentLoad = 0;

    for (int w : weights) {
        if (currentLoad + w > capacity) {
            days++;
            currentLoad = 0;
        }
        currentLoad += w;
    }
    return days;
}

int shipWithinDays(vector<int>& weights, int D) {
    int lo = *max_element(weights.begin(), weights.end());
    int hi = accumulate(weights.begin(), weights.end(), 0);

    while (lo <= hi) {
        int mid = lo + (hi - lo) / 2;
        if (daysNeeded(weights, mid) <= D) {
            hi = mid - 1;
        } else {
            lo = mid + 1;
        }
    }
    return lo;
}

int main() {
    vector<int> weights = {1,2,3,4,5,6,7,8,9,10};
    cout << "Minimum capacity: "
         << shipWithinDays(weights, 5) << endl;
    return 0;
}
\end{lstlisting}

\begin{complexitybox}
\textbf{Time Complexity.} The candidate capacity range halves each
iteration, costing $O(\log(\sum(\textit{weights})))$ iterations, and each
feasibility check costs $O(n)$ for the greedy simulation, giving
\[
  O(n \log(\sum(\textit{weights}))).
\]
\textbf{Space Complexity.} $O(1)$ auxiliary space.
\end{complexitybox}

\begin{takeawaybox}
When the items being processed must stay in a fixed order, the
feasibility check for binary search on the answer is usually a greedy
simulation (pack as much as possible into each ``bucket'' before starting
a new one) rather than a simple per-item summation. Verify, separately,
that the greedy simulation itself is optimal for the fixed-candidate
sub-problem -- here, via the exchange argument that under-loading a day
never helps -- before trusting it as the feasibility predicate.
\end{takeawaybox}
